\section{Literature Review}
\label{sec:lit}

\subsection{Classifying diabetes from survey data}
The study closest to mine is by Xie \emph{et al.}~\cite{xie2019}, who built type~2 diabetes risk models from the 2014 BRFSS. Their pipeline narrowed the data considerably: of 464,644 respondents, 138,146 remained after preprocessing, the outcome was restricted to non-pregnant adults over 30, and people with prediabetes were left out altogether. They compared eight classifiers on a two-thirds/one-third holdout and applied SMOTE only to the training portion, which is the correct way to resample. Discrimination was modest and tightly bunched, with AUCs between 0.718 and 0.795. The neural network scored best on AUC and accuracy yet caught only 37.8\% of cases, so the authors recommended the decision tree for screening because it found more of them. The paper taught me to judge screening models by sensitivity, and showed how much a filtered binary outcome can flatter results. One reported ``risk factor'' also stayed with me: a recent check-up was linked to \emph{higher} risk, almost certainly because people with diabetes see doctors more often. Survey predictors can reflect a diagnosis as much as anticipate it.

Most recent work uses a cleaned derivative of the 2015 BRFSS with 253,680 rows and 21 predictors, distributed through Kaggle and the UCI repository. Majcherek \emph{et al.}~\cite{majcherek2025} compared 18 algorithms on it and reported that an Extra Trees model reached an AUC of 0.99 and accuracy above 90\%, with BMI, age, general health and income as leading SHAP features. Zafar \emph{et al.}~\cite{zafar2025} applied ADASYN oversampling to the same data and proposed a hybrid deep architecture explained with LIME and SHAP. These results are hard to reconcile with the audit by Chowdhury \emph{et al.}~\cite{chowdhury2026}, which found AUROCs of only 0.818 to 0.830 across four learners, with gradient boosting already within 0.0003 of the estimated ceiling for what the survey variables can support. An AUC of 0.99 on data with such a ceiling suggests that something in the evaluation, perhaps oversampling or duplicate rows crossing the train--test boundary, deserves a closer look. I cannot tell the cause from the published papers, so I treat those figures as reported claims rather than benchmarks. The derivative dataset also makes two choices I wanted to avoid. Its binary version counts prediabetes as diabetes, which changes the meaning of the outcome, and it drops every respondent with a missing answer, even though people who refuse questions or stop interviews early are not a random subset.

\subsection{Simple models, complex models and tabular data}
Whether flexible models beat simple ones is an old question in clinical prediction, and the evidence is sobering. Christodoulou \emph{et al.}~\cite{christodoulou2019} reviewed 71 studies containing 282 comparisons and found no difference in discrimination between logistic regression and machine learning among comparisons at low risk of bias (difference in logit AUC 0.00, 95\% CI $-0.18$ to 0.18). The apparent advantages of machine learning were concentrated in the studies with weaker designs. Kopitar \emph{et al.}~\cite{kopitar2020} similarly found no clinically relevant gain from tree ensembles over regression for early type~2 diabetes detection.

The machine learning benchmarks tell a compatible story about neural networks. Grinsztajn \emph{et al.}~\cite{grinsztajn2022} tested deep architectures against tree ensembles on 45 datasets and found that trees remained state of the art on medium-sized tabular data. They traced the gap to inductive biases: trees cope well with uninformative features, respect the orientation of the feature axes and fit irregular functions easily. Shwartz-Ziv and Armon~\cite{shwartzziv2022} found that XGBoost beat recently proposed deep tabular models on most of the very datasets those models had been introduced with, and needed less tuning to do so. Whether this holds for health surveys, where nearly every predictor is binary or ordinal, is something I test directly. The ensembles I compare descend from Friedman's gradient boosting machine~\cite{friedman2001} and Breiman's random forests~\cite{breiman2001}, and my Functional API network adapts the wide-and-deep design of Cheng \emph{et al.}~\cite{cheng2016}, which places a linear model alongside a deep one.

\subsection{Evaluating rare outcomes}
Diabetes is a minority class in any general population survey, and prediabetes is rarer still. Saito and Rehmsmeier~\cite{saito2015} showed that ROC curves can look reassuring in this situation because the false-positive rate is computed against a large pool of negatives, while precision--recall curves, whose baseline equals the prevalence, show how hard the problem really is. That argument shaped my choice of average precision as the main metric. Resampling only helps when confined to training data, and published work is not always clear about this; the TRIPOD+AI statement~\cite{collins2024} now asks authors to report such choices along with missing-data handling and outcome definitions.

\subsection{Explanations, fairness and the quality of labels}
Impurity-based importances from random forests are biased towards variables with many categories or continuous scales~\cite{strobl2007}, so I rely instead on permutation importance and SHAP~\cite{lundberg2017}, which divides each prediction among the features using Shapley values. Fairness raises a deeper problem. Obermeyer \emph{et al.}~\cite{obermeyer2019} showed that a widely used population-health algorithm under-identified Black patients because it predicted health-care cost as a stand-in for health need. A similar trap exists here. A self-reported ``ever told'' diagnosis depends on access to care, so the groups who see doctors least are also the ones most likely to be labelled wrongly. Finally, a systematic review found that BRFSS self-reports agree reasonably with other self-report surveys but less well with surveys that take physical measurements~\cite{pierannunzi2013}, which limits how far any model trained on these labels can be read clinically.

\subsection{The gap this project addresses}
Taken together, these studies leave three gaps. Machine learning work on BRFSS diabetes data usually models either a filtered binary outcome~\cite{xie2019} or the complete-case derivative with prediabetes merged into diabetes~\cite{majcherek2025,zafar2025}, so the three categories the survey actually records are rarely modelled. Evaluations tend to stop at accuracy and ROC-AUC under default thresholds, saying little about minority-class precision, threshold choice or differences between groups of people. And comparisons of traditional and deep models seldom ask whether the extra complexity buys anything over a well-tuned linear model trained on the same, leakage-safe pipeline. My project works on the raw 2015 survey, keeps prediabetes as its own class, retains respondents with missing answers, and compares both paradigms under identical preprocessing with tuned decision rules and subgroup analysis. Individual pieces of this have been done before; the contribution lies in putting them together in one transparent and reproducible comparison.
